% ============================================================
\section{Foundational Proof Patterns}
\label{sec:foundational-lemmas}
% ============================================================

This section records the proof patterns used repeatedly in the clean horizon calculus.
They isolate the resolution split, the block decomposition, the telescoping identity, and the Pythagorean split that drive the later arguments.

% ------------------------------------------------------------
\subsection{Resolution-of-Identity Expansion}
\label{subsec:resolution-lemma}
% ------------------------------------------------------------

The basic split behind the calculus is the observable/shadow decomposition at a fixed horizon.

\begin{lemma}[Resolution-of-Identity Expansion]
\label{lem:resolution-identity}\lean{resolution_of_identity}
Let $\CCPi{h}$ be a horizon projection with shadow $\CCPibar{h} = I - \CCPi{h}$.
For any state $x$,
\[
x = \CCPi{h}x + \CCPibar{h}x.
\]
\end{lemma}

\begin{proof}
By definition, $\CCPibar{h}x = x - \CCPi{h}x$.
Hence
\[
\CCPi{h}x + \CCPibar{h}x = \CCPi{h}x + (x - \CCPi{h}x) = x.
\]
\end{proof}

\paragraph{Usage pattern.}
Whenever a proof says ``insert the horizon split,'' it means: replace a state
by its observable part plus its shadow part and then distribute.

% ------------------------------------------------------------
\subsection{Block Splitting Lemma}
\label{subsec:block-splitting}
% ------------------------------------------------------------

\begin{lemma}[Block Splitting for Operator Products]
\label{lem:block-splitting}\lean{block_splitting}
For additive operators $A, B$ and horizon projection $\CCPi{h}$,
\[
\CCRestr{(AB)}{h} = \CCObs{A}{h}\CCObs{B}{h} + \CCLeakCorr{h}{A}{B},
\]
where
\[
\CCLeakCorr{h}{A}{B} := \CCPi{h} A \CCPibar{h} B \CCPi{h}.
\]
\end{lemma}

(Proof in Appendix~\ref{sec:appendix-tower}.)

% ------------------------------------------------------------
\subsection{Three-Level Telescope}
\label{subsec:telescoping-lemma}
% ------------------------------------------------------------

\begin{lemma}[Three-Level Telescope]
\label{lem:orthogonal-telescoping}\lean{telescope_3_terms}
Let $(\CCPi{h})_{h \in \bbN}$ be a transitive horizon tower.
For horizons $h_0 \le h_1 \le h_2$ and any $x \in \cH$,
\[
(\CCPi{h_1} - \CCPi{h_0})x
+
\Bigl((\CCPi{h_2} - \CCPi{h_1})x + (x - \CCPi{h_2}x)\Bigr)
=
x - \CCPi{h_0}x.
\]
\end{lemma}

\begin{proof}
This is a direct telescoping regrouping:
\[
x - \CCPi{h_0}x
=
(\CCPi{h_1}x - \CCPi{h_0}x) + (\CCPi{h_2}x - \CCPi{h_1}x) + (x - \CCPi{h_2}x).
\]
The intermediate horizon contributions cancel.
\end{proof}

\paragraph{Calculus interpretation.}
This is the algebraic core of HFT-2: adjacent refinement layers close
internally, so only the outer remainder survives.

% ------------------------------------------------------------
\subsection{Pythagorean Norm Splitting}
\label{subsec:pythagorean-norm}
% ------------------------------------------------------------

\begin{lemma}[Pythagorean Norm Splitting]
\label{lem:pythagorean-splitting}\lean{pythagorean_splitting}\lean{pythagorean_splitting_observed}\lean{defectEnergy}
Assume the horizon tower is energy-orthogonal.
Then:
\begin{enumerate}[label=(\roman*)]
\item \textbf{State splitting:} for any $x \in \cH$,
\[
\|x\|^2 = \|\CCPi{h} x\|^2 + \|\CCPibar{h} x\|^2.
\]
\item \textbf{Operator action splitting:} for a closure operator $d$ and $x \in \CCHobs{h}$,
\[
\|dx\|^2 = \|\CCObs{d}{h} x\|^2 + \CCDefEnergy{h}{x},
\]
where
\[
\CCDefEnergy{h}{x} := \|\CCLeakExpr{d}{h}x\|^2.
\]
\end{enumerate}
\end{lemma}

\begin{proof}
(i) This is the exact energy decomposition at the horizon cut.

(ii) Apply (i) to $dx$.
If $x \in \CCHobs{h}$, then $\CCPi{h}x = x$, so the observable and shadow pieces of $dx$ are exactly $\CCObs{d}{h}x$ and $\CCLeakExpr{d}{h}x$.
\end{proof}

\begin{quote}\small
\textbf{Intuition.}
The horizon split does not create or destroy energy.
It partitions it into what is already visible and what remains in shadow.
\end{quote}
